\chapter{Introducción}
\label{chap1}
\ifpdf
  \graphicspath{{Chapter1/Chapter1Figs/PNG/}{Chapter1/Chapter1Figs/PDF/}{Chapter1/Chapter1Figs/}}
\else
  \graphicspath{{Chapter1/Chapter1Figs/EPS/}{Chapter1/Chapter1Figs/}}
\fi

\markboth{\hfill \thechapter. Introducción}{\hfill \thechapter. Introducción}

\section{Contexto y motivación}
\label{sec:contexto-motivacion}

% - epilepsia y EEG
La epilepsia es una enfermedad cerebral crónica no transmisible que afecta a personas de todas las edades y constituye uno de los trastornos neurológicos más frecuentes a nivel mundial \citep{WHO2024Epilepsy}. El electroencefalograma (EEG) es una herramienta de especial relevancia para apoyar el diagnóstico, la clasificación y el seguimiento de pacientes con epilepsia \citep{NoachtarRemi2009}. Mediante el registro de la actividad eléctrica cerebral, el EEG permite al especialista identificar patrones y eventos de interés para la evaluación neurológica. En la práctica clínica, estos estudios pueden generar registros extensos cuya revisión requiere tiempo, experiencia y herramientas adecuadas para facilitar su análisis.


% - contexto del Hospital Central del IPS
En el área de Neurología del Hospital Central del Instituto de Previsión Social (IPS), la adquisición y la revisión de las señales EEG se realizan mediante el \textit{software} Easy III, de Cadwell. Esta herramienta permite la captura y visualización de los registros electroencefalográficos. No obstante, a partir del relevamiento realizado en el contexto del proyecto se identificaron necesidades relacionadas con la integración del flujo de trabajo posterior a la adquisición, entre ellas el acceso mediante navegador web, la gestión centralizada de pacientes, estudios e informes clínicos y una mayor trazabilidad de las actividades realizadas durante el proceso de análisis.

% - flujo de trabajo actual
En el flujo de trabajo relevado, la especialista revisa los estudios en la misma computadora en la que está instalado el \textit{software} de adquisición, conectada al equipo de registro. Los informes se redactan en esa computadora mediante un procesador de textos, a partir de una plantilla, y se imprimen para su entrega al paciente. Tanto los estudios como los informes permanecen almacenados en ese equipo, por lo que su consulta depende del acceso a dicha computadora. A ello se suma la posibilidad de incorporar técnicas automatizadas de análisis de señales que permitan presentar información de apoyo al profesional médico, manteniendo la interpretación clínica y la toma de decisiones bajo la responsabilidad del especialista.

% - presentación de EpilexIA

En el marco del proyecto EpilexIA, el presente trabajo se centra en el desarrollo de una plataforma web concebida como \textit{software} de apoyo para el ámbito médico, orientada a complementar el flujo de trabajo existente en el área de Neurología del Hospital Central del IPS. La plataforma busca centralizar la gestión de pacientes, los estudios de EEG y los informes clínicos, permitir la visualización interactiva de señales electroencefalográficas y presentar las anomalías identificadas por los servicios de inteligencia artificial integrados al proyecto. Asimismo, incorpora mecanismos de control de acceso basados en roles y funciones de auditoría destinadas a favorecer la trazabilidad de las operaciones realizadas.
%CORREGIDO FIN:


% - marco institucional del proyecto
EpilexIA se desarrolla en el marco de un proyecto de investigación conjunta sustentado en una alianza entre tres sectores.
La cooperación internacional, a través de la Agencia de Cooperación Internacional de Corea (KOICA), aporta el financiamiento y el acompañamiento técnico del proyecto.
La academia, representada por la Facultad Politécnica de la Universidad Nacional de Asunción (FP-UNA) a través de su Centro de Innovación TIC (CIT), tiene a su cargo la investigación y el desarrollo del sistema y de los modelos de inteligencia artificial.
El sector privado, representado por la empresa BigQ, provee la infraestructura de cómputo para el entrenamiento de los modelos y el alojamiento del sistema en la nube.
El destinatario de la solución es el IPS, institución pública cuyo servicio de Neurología define los requerimientos, aporta los estudios y valida el sistema con sus especialistas.

Esta forma de colaboración permite orientar las capacidades de la cooperación internacional, la empresa y la universidad hacia un problema del sistema público de salud.
La participación del sector privado facilita el acceso a una infraestructura que de otro modo no estaría disponible para la investigación académica y favorece que sus resultados alcancen un entorno de uso real.

% - motivación para desarrollar EpilexIA
La motivación principal del proyecto radica en aprovechar tecnologías web para proporcionar una plataforma integrada y accesible desde diferentes dispositivos, como computadoras, \textit{notebooks} y \textit{tablets}, que complemente
 las herramientas utilizadas en el hospital. EpilexIA no pretende sustituir a Easy III en la etapa de adquisición del EEG, sino incorporar una capa complementaria para las etapas posteriores de carga, visualización, análisis asistido, revisión clínica, elaboración de informes y seguimiento del proceso.


El desarrollo de la propuesta consideró, además, las condiciones del
entorno en el que se llevó a cabo el proyecto, en particular la
disponibilidad de los profesionales del área para las instancias de
relevamiento y validación. En consecuencia, las funcionalidades de
EpilexIA se definieron y ajustaron progresivamente a partir de las
necesidades identificadas durante las actividades realizadas en el
Hospital Central del IPS, y se validaron en sesiones periódicas con una
médica especialista del servicio.




\section{Revisión de la literatura}
\label{sec:revision-literatura}

El electroencefalograma constituye una de las principales herramientas para el estudio de la actividad eléctrica cerebral y desempeña un papel relevante en la evaluación de pacientes con epilepsia. El análisis de los registros EEG permite identificar patrones asociados a diferentes síndromes epilépticos y aporta información útil para la localización y caracterización de la actividad epileptiforme. No obstante, su interpretación depende de diversos factores relacionados con el paciente, las condiciones del registro y los procedimientos utilizados durante el estudio \citep{NoachtarRemi2009}.

El desarrollo de herramientas computacionales ha permitido complementar el análisis convencional de las señales EEG mediante técnicas de procesamiento de señales y aprendizaje automático. Estas técnicas pueden emplearse para identificar características presentes en los registros electroencefalográficos y apoyar la detección de períodos asociados a crisis epilépticas u otros eventos de interés. En este ámbito, \citet{Shoeb2009} investigó la aplicación de métodos de aprendizaje automático para detectar el inicio de crisis epilépticas a partir de registros de EEG, lo que constituye un antecedente relevante para el desarrollo de métodos automatizados aplicados a este tipo de señales.

La investigación en este campo también se ha beneficiado de la disponibilidad de conjuntos de datos de EEG utilizados para el desarrollo y la evaluación de métodos computacionales. Un ejemplo es la CHB-MIT Scalp EEG Database, publicada en PhysioNet, que contiene registros de pacientes pediátricos con crisis epilépticas y anotaciones asociadas a dichos eventos \citep{Guttag2010CHBMIT}. Estos recursos han favorecido el estudio de métodos orientados al procesamiento y a la detección automatizada de eventos en señales de EEG.

Estos antecedentes evidencian la incorporación progresiva de herramientas computacionales como apoyo al procesamiento y al análisis de registros electroencefalográficos. Sin embargo, la identificación automatizada de eventos constituye solo una parte del proceso de análisis. Para su utilización en un entorno clínico, los resultados obtenidos deben presentarse de manera que el profesional pueda acceder al estudio, visualizar las señales, revisar la información generada por los métodos automatizados y utilizarla como elemento de apoyo durante la evaluación y la elaboración del informe clínico.



Esta evolución se relaciona también con el desarrollo de sistemas de
\textit{software} destinados a cumplir funciones específicas dentro del ámbito
médico. En este contexto adquiere relevancia el concepto de \textit{software}
como dispositivo médico (\textit{Software as a Medical Device}, SaMD),
definido por el Foro Internacional de Reguladores de Dispositivos
Médicos (IMDRF) \citep{IMDRF2013SaMD}, así como las consideraciones
asociadas al uso de herramientas informáticas que proporcionan
información de apoyo con fines clínicos. Estos conceptos, junto con
los fundamentos de la electroencefalografía, el procesamiento de
señales y la inteligencia artificial aplicada al EEG, se desarrollan
con mayor profundidad en el marco teórico y en el estado del arte
presentados en el Capítulo~2.


Dentro de este contexto, EpilexIA se plantea como una plataforma web que reúne la gestión de pacientes y estudios de EEG, la visualización de señales electroencefalográficas, la presentación de anomalías registradas a partir de un servicio de inteligencia artificial integrado y la elaboración de informes clínicos. De esta manera, el proyecto no se limita al análisis automatizado de las señales, sino que busca incorporar sus resultados en un entorno que apoye el flujo de trabajo del profesional especialista.


\section{Objetivos}
\label{sec:objetivos}

\subsection{Objetivo general}
\label{subsec:objetivo-general}


Desarrollar un sistema web de apoyo a la detección de crisis epilépticas
mediante el análisis de señales electroencefalográficas (EEG) asistido
por inteligencia artificial, que permita gestionar pacientes, estudios
EEG e informes clínicos, así como visualizar las anomalías identificadas
por el componente de inteligencia artificial integrado, con el propósito
de complementar el flujo de trabajo existente con Easy III en el
área de Neurología del Hospital Central del Instituto de Previsión Social.




\subsection{Objetivos específicos}
\label{subsec:objetivos-especificos}

\begin{enumerate}
    \item Relevar el estado del arte sobre sistemas de información clínica para la epilepsia, la visualización de señales de EEG en entornos web y las herramientas de apoyo al diagnóstico neurológico.

    \item Analizar los requerimientos funcionales y no funcionales del área de Neurología del Hospital Central del IPS, identificando las limitaciones del flujo actual con Easy III que el sistema web debe resolver.

    \item Diseñar e implementar una base de datos relacional en PostgreSQL que soporte la gestión de pacientes, estudios de EEG, archivos de señal, informes clínicos, anomalías detectadas y auditoría de actividad.

    \item Desarrollar un \textit{backend} REST (Java/Spring Boot) con control de acceso basado en roles para especialistas, técnicos y administradores, que gestione los datos clínicos y se integre con el servicio de inteligencia artificial mediante su API.

    \item Construir un \textit{frontend} web adaptable (\textit{responsive}) con React y TypeScript, que incluya un visualizador interactivo de señales EEG, un módulo de gestión de anomalías provenientes del servicio de IA, un editor de informes clínicos y un panel de auditoría, accesible desde computadoras, \textit{notebooks} y \textit{tablets}.


    \item Evaluar el sistema mediante pruebas funcionales, de usabilidad e integración con registros EEG reales del Hospital Central del IPS.
\end{enumerate}


\section{Problema y oportunidad}
\label{sec:problema-oportunidad}

En el área de Neurología del Hospital Central del Instituto de Previsión Social, la adquisición y revisión de los estudios electroencefalográficos se realizan con el Easy III. A partir del relevamiento realizado en el marco del proyecto, se identificaron necesidades relacionadas con el acceso a los estudios, la integración de la información y la trazabilidad de las actividades que forman parte del flujo de trabajo posterior a la adquisición de las señales EEG.

Una de las necesidades identificadas es el acceso a los registros electroencefalográficos desde distintos dispositivos. La dependencia de equipos con el \textit{software} instalado limita la
posibilidad de consultar los estudios en un navegador web desde computadoras, \textit{notebooks} o \textit{tablets}, por lo que el acceso a la información queda condicionado a la disponibilidad de los equipos y herramientas utilizados en el flujo de trabajo existente.


Asimismo, la información relacionada con pacientes, estudios EEG e informes clínicos puede encontrarse distribuida entre las herramientas utilizadas en las diferentes etapas del proceso.
 Esta situación dificulta contar con un entorno unificado que reúna la información necesaria para la revisión del estudio y el seguimiento de las actividades realizadas al respecto.

Otra necesidad se refiere a la trazabilidad de las operaciones y al control de acceso. El manejo de información asociada a pacientes, estudios e informes requiere mecanismos que permitan registrar las acciones efectuadas por los usuarios y establecer permisos de acuerdo con las responsabilidades de los distintos perfiles que intervienen en el proceso.

A estas necesidades se suma la incorporación de herramientas computacionales
que apoyen la revisión de registros EEG extensos. La identificación
automatizada de posibles anomalías puede proporcionar información
complementaria al especialista y facilitar la localización de segmentos de
interés dentro del registro. Estos resultados, sin embargo, deben utilizarse
como elementos de apoyo y permanecer sujetos a la revisión e interpretación
del profesional especialista.


Las necesidades identificadas representan una oportunidad para
complementar el flujo de trabajo existente mediante una plataforma web
que centralice la información asociada a los estudios EEG y facilite
su acceso desde diferentes dispositivos. En este contexto, EpilexIA
propone integrar en un mismo entorno la gestión de pacientes y estudios,
la visualización de señales EEG, la presentación de resultados
provenientes del procesamiento mediante inteligencia artificial,
la elaboración de informes clínicos, el control de acceso y la
auditoría de las actividades realizadas. De esta manera, se plantea
una solución unificada de apoyo al trabajo de los profesionales del
área de Neurología.


\section{Justificación}
\label{sec:justificacion}

El desarrollo de EpilexIA se justifica por la necesidad de disponer de una
herramienta que complemente el proceso de análisis de estudios
electroencefalográficos en el área de Neurología del Hospital Central del
Instituto de Previsión Social. La utilización de una plataforma web permite integrar en un mismo
entorno diferentes actividades asociadas con la revisión de los
estudios EEG y facilitar el acceso a la información necesaria
durante el trabajo de los profesionales del área.


Desde el punto de vista operativo, la centralización de pacientes, estudios
EEG e informes clínicos permite organizar y relacionar la información
correspondiente a las distintas etapas del proceso. A ello se suman
mecanismos de control de acceso y auditoría, mediante los cuales es posible
establecer permisos de acuerdo con las responsabilidades de cada perfil de
usuario y registrar las operaciones realizadas dentro del sistema. Estas
características están orientadas a favorecer la trazabilidad de las
actividades y el seguimiento de los estudios.

El enfoque web permite además ampliar las posibilidades de acceso a las
funcionalidades de la plataforma. Mediante tecnologías compatibles con
navegadores modernos y una interfaz adaptable, EpilexIA puede utilizarse
desde computadoras, \textit{notebooks} y \textit{tablets} sin requerir
la instalación de la aplicación en cada dispositivo. Esta característica proporciona una alternativa para la consulta y revisión de la información asociada con los
estudios EEG desde los dispositivos compatibles con el sistema.

Otro elemento que fundamenta el proyecto es la integración de un servicio
de inteligencia artificial para el análisis de señales EEG. Los resultados
generados por este componente pueden presentarse dentro del sistema como
información complementaria para apoyar la identificación y revisión de
segmentos de interés. La incorporación de esta capacidad no pretende
sustituir el criterio del especialista ni establecer de manera autónoma un
diagnóstico clínico, sino proporcionar información de apoyo que pueda ser
evaluada por el profesional junto con las señales EEG y los demás datos
disponibles en el estudio.

Desde una perspectiva tecnológica y académica, EpilexIA permite aplicar e
integrar tecnologías de desarrollo web, gestión de bases de datos,
procesamiento de señales e inteligencia artificial en un problema del ámbito
de la salud. El proyecto comprende no solamente la implementación de la
plataforma, sino también el relevamiento de requerimientos, el diseño
de su arquitectura, la integración de sus diferentes componentes
y la realización de pruebas funcionales del sistema. De esta manera, se busca desarrollar una solución tecnológica acorde con las necesidades identificadas en el área de Neurología del Hospital Central del IPS y generar conocimiento y experiencia aplicables al desarrollo de sistemas informáticos
de apoyo al análisis de información clínica.


\section{Solución propuesta}
\label{sec:solucion-propuesta}

Como respuesta a las necesidades identificadas, se propone EpilexIA, una
plataforma web orientada a complementar el flujo de trabajo relacionado con
el análisis de estudios electroencefalográficos en el área de Neurología del
Hospital Central del Instituto de Previsión Social. El sistema busca
centralizar en un único entorno las principales actividades posteriores a la
adquisición del EEG, proporcionando herramientas para la gestión de pacientes
y estudios, la visualización de señales, la revisión de anomalías y la
elaboración de informes clínicos.

La plataforma contempla la gestión de la información necesaria para asociar
cada estudio EEG con el paciente correspondiente y mantener organizados los
datos generados durante su procesamiento y revisión. Los archivos de señales
pueden ser incorporados al sistema para su posterior visualización mediante
un visor interactivo, desde el cual el profesional especialista puede explorar
los diferentes canales del registro y utilizar herramientas destinadas a
facilitar su análisis.


EpilexIA incorpora además la integración con un servicio de inteligencia
artificial destinado al procesamiento automatizado de señales EEG.
La plataforma está diseñada para incorporar los resultados de este
componente y presentarlos al profesional como información de apoyo
para la identificación y revisión de segmentos de interés dentro
del registro. Estos resultados no constituyen por sí mismos un
diagnóstico clínico y su interpretación permanece bajo la
responsabilidad del profesional especialista.


Como parte del flujo de revisión, el sistema permite elaborar y gestionar
informes clínicos asociados con los estudios EEG. De esta manera, la visualización de las señales, la consulta de anomalías registradas en la plataforma y la elaboración del informe
pueden realizarse dentro de un mismo entorno, manteniendo la relación entre la información correspondiente a cada estudio.


La solución contempla asimismo mecanismos de autenticación, autorización y
control de acceso basados en los perfiles definidos para los usuarios del
sistema. Las operaciones relevantes son registradas mediante funcionalidades
de auditoría, con el propósito de proporcionar trazabilidad sobre las acciones
realizadas dentro de la plataforma y facilitar el seguimiento de las
actividades relacionadas con los estudios.

Desde el punto de vista tecnológico, EpilexIA se estructura como una
aplicación web compuesta por una interfaz de usuario, servicios de \textit{backend},
una base de datos relacional y componentes destinados a la integración y
ejecución del procesamiento mediante inteligencia artificial. Esta separación
permite distribuir las responsabilidades entre los diferentes componentes
del sistema y favorecer su mantenimiento y evolución. La arquitectura y las
decisiones de diseño se presentan en el Capítulo~4, mientras que los aspectos
relacionados con su implementación e integración se desarrollan en el
Capítulo~5.

En conjunto, la solución propuesta busca proporcionar un entorno integrado
que complemente las herramientas existentes en el Hospital Central del IPS.
EpilexIA no sustituye al sistema utilizado para la adquisición de las
señales EEG, sino que complementa las etapas posteriores del flujo de
trabajo mediante la organización de la información, la visualización
de señales, el apoyo al análisis, la revisión clínica, la elaboración
de informes y la trazabilidad de los estudios.



\section{Alcance y delimitaciones}
\label{sec:alcance}

El alcance del proyecto comprende el desarrollo de una plataforma web
orientada a apoyar el análisis y la gestión de estudios EEG en el área de
Neurología del Hospital Central del Instituto de Previsión Social. El sistema
abarca la información necesaria para identificar a los pacientes y asociarlos
con sus respectivos estudios, la incorporación y gestión de archivos de
señales EEG, su visualización interactiva, la presentación de resultados
provenientes del procesamiento mediante inteligencia artificial y la
elaboración de informes clínicos asociados con los estudios.

La plataforma contempla además mecanismos de autenticación y autorización
para controlar el acceso a las diferentes funcionalidades según el perfil
del usuario, así como el registro de operaciones mediante funciones de
auditoría. El acceso se realiza mediante una interfaz web adaptable a
diferentes tamaños de pantalla, con el propósito de permitir la utilización
del sistema desde computadoras, \textit{notebooks} y \textit{tablets}
compatibles con navegadores web modernos.


El proyecto se delimita a las etapas posteriores a la adquisición de las
señales electroencefalográficas. Por tanto, EpilexIA no controla los equipos
utilizados para realizar el electroencefalograma ni sustituye al \textit{software}
empleado para la adquisición original de los registros. Los archivos
generados durante este proceso constituyen la información de entrada que
posteriormente puede ser incorporada y procesada dentro de la plataforma.


EpilexIA contempla la integración de un componente de inteligencia
artificial destinado al procesamiento automatizado de señales EEG.
El alcance del presente trabajo comprende la integración de dicho
componente con los demás elementos del sistema y la presentación de
sus resultados al profesional especialista. El desarrollo y
entrenamiento del modelo de inteligencia artificial no forman parte
del alcance de este trabajo.

Las verificaciones documentadas permitieron comprobar el funcionamiento integral del procesamiento automatizado sobre un conjunto reducido de estudios reales.
La evaluación de su eficacia diagnóstica a mayor escala excede el alcance de este trabajo. Asimismo, los
resultados obtenidos mediante inteligencia artificial constituyen
información de apoyo y no representan por sí mismos un diagnóstico
médico. La evaluación e interpretación clínica permanecen bajo la
responsabilidad del profesional especialista.


La gestión de pacientes incluida en el sistema se limita a la información
necesaria para su identificación y para la asociación con los estudios EEG
gestionados por la plataforma. En consecuencia, EpilexIA no pretende
reemplazar los sistemas institucionales destinados a la administración
integral de historias clínicas ni centralizar la totalidad de la información
médica del paciente.


Finalmente, las pruebas y verificaciones realizadas se circunscriben al contexto y a los requerimientos relevados para el área de Neurología del Hospital Central del IPS. Si bien la arquitectura y las tecnologías utilizadas permiten considerar futuras adaptaciones o ampliaciones, la utilización del sistema en otros entornos clínicos podría requerir nuevos relevamientos, ajustes funcionales
y procesos de validación específicos.

\section{Hipótesis del proyecto}
\label{sec:hipotesis}

La implementación de un sistema web integrado que centralice la gestión de
pacientes, estudios EEG e informes clínicos, e incorpore la visualización de
resultados provenientes del procesamiento mediante inteligencia artificial,
puede contribuir a mejorar el flujo de trabajo relacionado con el análisis
de estudios electroencefalográficos en el área de Neurología del Hospital
Central del Instituto de Previsión Social.

En particular, se plantea que la utilización de EpilexIA puede facilitar el
acceso a los estudios desde diferentes dispositivos, favorecer la
trazabilidad de las actividades realizadas durante su revisión y proporcionar
herramientas de apoyo para la localización y evaluación de segmentos de
interés dentro de registros EEG extensos. La integración de estas
funcionalidades en un mismo entorno podría contribuir además a facilitar
determinadas etapas del proceso de revisión y elaboración de informes
clínicos, complementando el flujo de trabajo existente.


La hipótesis considera el procesamiento mediante inteligencia artificial como
un mecanismo de apoyo al análisis y no como un sustituto de la evaluación
médica. En consecuencia, los posibles beneficios derivados del uso del
sistema se relacionan con la organización, accesibilidad, trazabilidad y
asistencia durante el proceso de revisión, mientras que la interpretación de
los registros EEG y las decisiones clínicas continúan siendo responsabilidad
del profesional especialista.


\section{Resultados esperados}
\label{sec:resultados-esperados}

Como resultado del desarrollo del proyecto se espera disponer de una
plataforma web funcional que permita integrar las principales actividades
relacionadas con la gestión y el análisis de estudios EEG en el área de
Neurología del Hospital Central del Instituto de Previsión Social. La
solución deberá proporcionar un entorno accesible mediante navegador web
y adaptable a los diferentes dispositivos contemplados dentro del alcance
del proyecto.

Se espera que el sistema permita gestionar la información necesaria de los
pacientes y asociarla con sus respectivos estudios EEG e informes clínicos,
manteniendo organizados los datos generados durante las diferentes etapas
del proceso. Asimismo, se espera disponer de mecanismos de autenticación,
control de acceso y auditoría que permitan gestionar las responsabilidades
de los usuarios y proporcionar trazabilidad sobre las operaciones realizadas
en la plataforma.


Otro resultado esperado corresponde a la disponibilidad de un visualizador
interactivo de señales electroencefalográficas que proporcione al profesional
especialista herramientas para explorar y revisar los registros EEG. Este
visualizador deberá permitir la consulta de las anomalías identificadas
mediante el componente de inteligencia artificial, de manera que puedan
ser localizadas y examinadas por el especialista en conjunto con las
señales EEG dentro del mismo entorno de trabajo.


Se espera además que la plataforma permita elaborar y gestionar informes
clínicos asociados con los estudios EEG, integrando esta actividad con la
visualización de las señales y la revisión de los resultados del análisis
automatizado. Esta integración busca proporcionar continuidad entre las
diferentes etapas posteriores a la adquisición del estudio y mantener
relacionada la información utilizada durante su evaluación.

Como parte de los resultados del proyecto, se prevé también disponer de la
documentación técnica y de usuario necesaria para describir el funcionamiento
del sistema, sus principales componentes y los procedimientos requeridos
para su utilización. Esta documentación deberá servir de apoyo tanto para
los usuarios de la plataforma como para las tareas posteriores de
mantenimiento y evolución del sistema.


Finalmente, se espera evaluar el funcionamiento de EpilexIA mediante
pruebas funcionales, de integración y de usabilidad, considerando
registros EEG y el contexto de utilización definido para el Hospital
Central del IPS. Estas pruebas permitirán verificar el cumplimiento
de los requerimientos funcionales establecidos y analizar los
resultados obtenidos respecto de los objetivos planteados para el
proyecto, sin que ello implique una validación de la eficacia
diagnóstica del componente de inteligencia artificial.



\section{Metodología general}
\label{sec:metodologia-general}

Para el desarrollo de EpilexIA se adoptó el Modelo V como metodología
general de trabajo. Este enfoque permite establecer una correspondencia
entre las etapas de definición y diseño del sistema y las actividades
posteriores de verificación y validación, de manera que los requerimientos
y decisiones establecidos durante el desarrollo puedan ser contrastados
mediante los niveles de prueba correspondientes. Esta relación entre
las actividades de desarrollo y las instancias de verificación y
validación se encuentra fundamentada en el enfoque de ingeniería de
sistemas descrito por \citet{ForsbergMooz1991}.


En el proyecto EpilexIA, la aplicación de esta metodología parte del
relevamiento de requerimientos realizado con profesionales del área de
Neurología del Hospital Central del Instituto de Previsión Social. 
A partir de las necesidades identificadas se realizaron las actividades
de análisis y diseño del sistema, seguidas por la implementación e
integración de sus componentes. Posteriormente, se llevaron a cabo pruebas orientadas a verificar
el comportamiento de las funcionalidades desarrolladas, su integración
y su correspondencia con los requerimientos establecidos.



La aplicación del Modelo V fue adaptada a las condiciones particulares del
proyecto y a las posibilidades de interacción con los profesionales del
Hospital Central del IPS. Las instancias de trabajo con los usuarios se utilizaron para relevar
necesidades, presentar avances y obtener observaciones sobre las
funcionalidades desarrolladas. La retroalimentación obtenida permitió
realizar ajustes durante el desarrollo y mantener la correspondencia
entre los requerimientos identificados, la solución implementada y
las actividades de verificación realizadas.


Esta adaptación tuvo en cuenta la disponibilidad de los profesionales
del área de Neurología para participar en las actividades del proyecto.
Por este motivo, las instancias de interacción se organizaron de
acuerdo con sus posibilidades de participación, procurando aprovechar
cada encuentro para revisar los requerimientos, evaluar los avances
del sistema y recopilar observaciones que contribuyeran a su mejora.


La descripción detallada del Modelo V, su adaptación al proyecto EpilexIA,
las actividades realizadas y los mecanismos de verificación y validación
empleados se presenta en el Capítulo~3.


\section{Organización del trabajo}
\label{sec:organizacion}

El presente trabajo se encuentra organizado en siete capítulos, estructurados
de manera que permitan presentar progresivamente los fundamentos, el proceso
de desarrollo, el diseño, la implementación y la evaluación de EpilexIA.

El Capítulo 1 presenta la introducción al proyecto, incluyendo el contexto y
la motivación, la revisión de la literatura, los objetivos, el problema y la
oportunidad identificados, la justificación, la solución propuesta, el alcance
y las delimitaciones, la hipótesis, los resultados esperados y la metodología
general adoptada.

El Capítulo 2 desarrolla el marco teórico y el estado del arte necesarios para
comprender los fundamentos del proyecto. Se abordan el concepto de \textit{software}
como dispositivo médico (SaMD), la epilepsia, la electroencefalografía, la
representación y el procesamiento de señales EEG y la aplicación de técnicas
de inteligencia artificial a su análisis. Asimismo, se presentan sistemas y
herramientas relacionados con la gestión y visualización de EEG, junto con
antecedentes nacionales e internacionales relevantes para el desarrollo de
EpilexIA.

El Capítulo 3 describe la metodología de desarrollo utilizada durante el
proyecto, detallando el Modelo V y su adaptación al contexto de EpilexIA.
Se presentan las actividades de relevamiento de requerimientos, diseño,
desarrollo, verificación, integración y validación, así como los mecanismos
de retroalimentación empleados durante el proceso.

El Capítulo 4 presenta el análisis de requerimientos y el diseño de EpilexIA.
En este capítulo se describen el contexto operativo, los actores y perfiles de
usuario, los requerimientos funcionales y no funcionales y los diagramas de
secuencia. Asimismo, se desarrolla el diseño de la arquitectura, la base de
datos, los servicios e interfaces de integración, el componente de inteligencia
artificial integrado en EpilexIA, los mecanismos de seguridad, auditoría y
trazabilidad, y la interfaz de usuario.

El Capítulo 5 describe la implementación de EpilexIA, incluyendo las
tecnologías y herramientas utilizadas y la construcción e integración de los
principales componentes de la plataforma. Se presentan la implementación de la
base de datos, el \textit{backend} y el \textit{frontend}, la gestión de pacientes y estudios
EEG, el procesamiento y visualización de las señales, la integración con el
servicio de inteligencia artificial, la gestión de anomalías e informes
clínicos, el control de acceso, la auditoría y el despliegue de los componentes
del sistema.

El Capítulo 6 presenta las pruebas y verificaciones realizadas sobre
el sistema. Se describen las pruebas funcionales y de integración,
así como las verificaciones relacionadas con la visualización de
señales EEG y la presentación de anomalías registradas en la
plataforma. Asimismo, se exponen los resultados de las actividades
de evaluación realizadas, sus principales limitaciones y el grado
de cumplimiento de los requerimientos establecidos, considerando
los objetivos planteados para el proyecto.



Finalmente, el Capítulo 7 presenta las conclusiones derivadas del
desarrollo y la evaluación de EpilexIA, analiza el grado de
cumplimiento de los objetivos y la evidencia obtenida en relación
con la hipótesis planteada, identifica los principales aportes
y limitaciones del proyecto y establece posibles líneas de
trabajo futuro.
